Formal Analysis of BPMN Models Using
Event-B
Jeremy W. Bryans1 and Wei Wei2
1 School of Computing Science, Newcastle University, United Kingdom
Jeremy.Bryans@ncl.ac.uk 2 SAP Research CEC Darmstadt, SAP AG, Bleichstr. 8, 64283 Darmstadt, Germany
wei01.wei@sap.com
Abstract. The use of business process models has gone far beyond documentation purposes. In the development of business applications, they
can play the role of an artifact on which high level properties can be
verified and design errors can be revealed in an effort to reduce overhead
at later software development and diagnosis stages. This paper demonstrates how formal verification may add value to the specification, design and development of business process models in an industrial setting.
The analysis of these models is achieved via an algorithmic translation
from the de-facto standard business process modeling language BPMN
to Event-B, a widely used formal language supported by the Rodin platform which offers a range of simulation and verification technologies.
Keywords: business process modelling, verification, BPMN, Event-B.
1 Introduction
Complex, large-scale business information systems are critical to the successful
operation of many businesses, and SAP is a leading provider of such systems.
Business process modeling has become increasingly important to the development of enterprise software applications [13]. Nowadays, business applications
are usually built by integrating a broad range of highly configurable software
components and services, which can be rapidly tailored to satisfy different and
constantly changing business needs. Business process models are used to describe such integration scenarios and their work flows, facilitating an intuitive
common understanding of the business logic between customers and developers. In addition to their use as documentation, business process models can also
be simulated, analyzed and verified to reveal design errors at an early stage in
software development. This promises to enhance the efficiency of reaching highquality software solutions and can save substantial implementation and diagnosis
costs which would otherwise be incurred at later development phases.
We wish to use formal methods to improve the quality of business process
models within a software design process, and also aim to reduce the extra burden that formal methods induce on designers and developers. Within the context
of the DEPLOY project1, we choose the Event-B modeling formalism [1] and the
1 www.deploy-project.eu
S. Kowalewski and M. Roveri (Eds.): FMICS 2010, LNCS 6371, pp. 33–49, 2010.
c Springer-Verlag Berlin Heidelberg 2010
34 J.W. Bryans and W. Wei
Rodin platform [2] in our pursuit of these goals. The choice is also encouraged
by our past successful experiences of using Event-B for describing and analyzing business applications [5,6]. Event-B offers many indispensable features for
analyzing business process models such as the ability to model data. The Rodin
platform is empowered by a large number of plug-ins providing various analysis
capabilities like specialized provers, model checking, and simulation.
This paper examines our recent work on the formal analysis of business process models using Event-B and Rodin, and discusses the impact of the analysis
results on software design and development. We also investigate the potential to
largely automate these analyses in order to pave the way for future industrial
deployment. We designed an algorithmic translation from BPMN, the de-facto
standard business process modeling language, to Event-B. The translation covers most of the commonly used BPMN features, also including features newly
introduced in the proposed draft of the second version of the language [15]. We
also make the Event-B translation structurally faithful to the original BPMN
model, which not only improves readability, but also enhances provability and
analyzability.
Outline. In Section 2 we briefly introduce BPMN and Event-B, and in Section 3
sketch our translation from BPMN to Event-B. Sections 4 and 5 describe two
case studies to illustrate how formal analysis is performed on the Event-B translations of BPMN models, and also discuss the possibility of automating these
analysis procedures. Related work is discussed in Section 6, before we conclude
in Section 7 with a discussion of our next steps. Due to space constraints, we
have moved some of the Event-B code and discussion into appendices.
2 Background
BPMN. We introduce the Business Process Modeling Notation (BPNM) elements we use in this paper. We only show syntactic compositions here. The
semantics of syntactic elements will be discussed later when we explain how
they are translated.
A typical BPMN model consists of one or more pools, each representing a collaboration partner (such as FACTORY and WORKER in Figure 1). Each pool usually
contains a top process. A process contains flow objects and the connections between them. Flow objects include events, gateways and activities. Events either
throw or catch triggers and are represented as circles containing a marker indicating the kind of trigger. Gateways converge or diverge control flows and are
represented as diamonds. An activity can be either an atomic task or a composite
sub-process that contains an inner process. An activity can be a loop. Activities
are graphically represented as rounded rectangles. A process may contain data
items as process instance attributes. There are also data stores that are processindependent and globally accessible.
Two pools communicate with each other mainly by exchanging messages,
which may carry data fields. Message flows are represented as directed dotted
Formal Analysis of BPMN Models Using Event-B 35
lines connecting two pools. BPMN does not dictate how the message exchange
mechanism works. In this paper, we assume that messages may not be lost,
duplicated, or altered but may arrive in any order. Furthermore, BPMN provides
the concept of correlation to identify the proper recipient of a message.
An important concept of BPMN is activity compensation that usually happens when the effect of an activity is no longer desired and needs to be reversed.
We will discuss compensation in greater length in Section 3.6. A complete description of all BPMN features can be found in [15].
Event-B and Rodin. An Event-B model consists of contexts and machines. The
contexts describe the static elements of the model, whereas the machines specify
the dynamic behavior of the model. Each machine may contain variables that
model persistent state data, invariants that restrict the valid content of variables, and guarded events that describe functionality of the machine in terms of
actions defined over the state variables. Typically, a model consists of a chain of
Event-B machines, each of which (apart from the first) is linked to its predecessor by a refinement relation expressed in terms of a gluing invariant between the
two machines. In a refinement relation, we refer to the successor machine and its
components as concrete and the predecessor and its components as abstract.A
concrete event refines an abstract event if the guards of the concrete event imply
the guards of the abstract event and the abstract actions simulate the concrete
ones with respect to the gluing invariant. Machines and refinement steps give
rise to proof obligations that ensure internal consistency of individual machines
(e.g. well-definedness and feasibility of events) and behavior preservation across
refinement steps. A typical Event-B model has an extremely simple initial machine, with detail added in a controlled way through refinement steps. These
steps are usually small to reduce the size and complexity of the generated proof
obligations and the associated burden on the automatic provers. We make substantial use of refinement in our translation from BPMN to Event-B. A detailed
account of the Event-B language can be found in [1].
Rodin is an open, extensible toolset for modelling and verification of Event-B
models. A model editing interface is provided for constructing Event-B machines
and refinement steps. Proof obligations are automatically generated and discharged (as far as possible automatically) by proof tools built into the platform.
In the event of an obligation not being automatically discharged, an interface
for manual proof guidance can used.
3 Translating BPMN to Event-B
BPMN is specified using natural and graphic languages, and comes with no
rigorous semantics defined. Therefore, there are a lot of ambiguities in BPMN
that had to be clarified when we designed the translation into Event-B. These
clarifications are according to the specific needs of our use cases, so by no means
do they offer the only proper solutions – other semantic variants can be chosen.
The translation covers most of the commonly used BPMN features including
comprehensive modeling of control flows, data modeling, compensation, message
36 J.W. Bryans and W. Wei
based communication, error and exception handling, sub-processes, looping and
multi-instance activities. The uncovered BPMN features are most notably choreography and conversations as well as some types of flow objects, including call
activities, transactions, conditional events and complex gateways. Some of these
missing features are rarely used in practice and add significant complexity to the
model. Other missing features such as transactions have very vague descriptions
in the official BPMN specification and are difficult to interpret.
Our translation was guided by three principles. First, the Event-B translation
should be structurally faithful to the original BPMN model so that anyone
with knowledge of the original model can easily understand the translation.
Also, any analysis result that we may obtain from the Event-B translation can
be easily mapped back to the original model. Second, the translation should be
designed to improve provability, i.e. it should result in the automatic discharge
of as many proof obligations as possible. Finally, we are interested in verifying
properties for systems that allow multiple instances of same processes.
We are unable to give a detailed description of how each BPMN element is
translated. We therefore select a few important BPMN features and explain the
intuition of their translation.
3.1 The Structure of the Translation
We take the model in Figure 1 as an example to show how its Event-B translation
is structured (Figure 2). This model describes the management of shift work
within a factory: A worker assigned to a shift becomes unavailable, and the
manager has to find a replacement from the pool of available workers. The status
of each worker is maintained in a database. In this scenario, an attempt is made
to automatically choose a replacement. A request is sent to an available worker,
who has a fixed length of time to reply. If he accepts, he is assigned to the shift
and the database is updated. Otherwise, the process may be repeated up to a
maximum of five times. If, after five attempts, a replacement has not been found,
a manager steps in to allocate a worker to the shift directly.
The contexts in the Event-B translation contain common definitions such as
process life cycle states as well as abstract constants and carrier sets that represent process instances, message instances, data types, and so on. The translations
of processes and their communication are gradually added to a series of refining
machines: The machine at the first level contains nothing but the control flow
information of the Factory process. In particular, it has neither data information nor the internal detail of the sub-process schedule. The machine at the
second level preserves or refines all information in the first machine, and adds
also the data flow information of Factory. Details of schedule and WORKER are
added similarly into later refinements. In the end, the communication between
the two top processes is added into the last machine.
The above structure preserves the hierarchical structure of the original model
through refinements: the information of a sub-process (e.g., schedule) is always
added into machines at higher refinement levels than that of the container process
(e.g., FACTORY). Our structure also achieves separation of concerns, which is very
Formal Analysis of BPMN Models Using Event-B 37
Fig. 1. The shift worker scheduling model
Fig. 2. The structure of the Event-B translation of Figure 1
beneficial for automated provers: A property about the control flow of process
Factory can be expressed and proved at the first refinement level since it needs
no information from later levels. This means a smaller hypothesis space for
automated provers to search.
3.2 Processes
We allow multiple instances of a process. We use an abstract carrier set to
represent all possible instances of each process (e.g. PROC FACTORY INSTANCES)
in contexts. The machines contain variables recording existing process instances
(e.g. instances Factory); recording the life cycle state of each existing instance
(e.g. state Factory) and, in case of a sub-process, recording the parent of each
sub-process instance (e.g. parent inner schedule); and recording which activity instance (outer instance) results in the creation of a sub-process instance
(e.g. outer inner schedule).
38 J.W. Bryans and W. Wei
Control flow. Our interpretation of sequential and parallel executions of flow
objects uses tokens. For each sequence flow, we define a function that maps each
process instance to the number of tokens in this particular process instance.
Tokens are initialized when a new process instance is created by a start event:
all outgoing sequence flows from the start event receive a certain number of
tokens (usually 1), and all other flows receive no tokens. Each flow object is
guarded by a condition stipulating how many tokens it needs to start execution.
Control flow convergence. With a few exceptions like join gateways, a flow object with multiple incoming flows needs only one of the incoming flows to carry
enough tokens to start execution. In this case, we use as many Event-B events
to represent the flow object as the number of incoming flows: each event describes the situation in which the tokens on the corresponding incoming flow are
consumed. This is because otherwise we must express disjunctive choices and
updates of tokens in the guard and action of the Event-B event representing the
flow object. Automated provers often struggle to deal with disjunctions because
they lead to case splitting and a potential explosion in the size of the proof tree.
A example of the translation of control flow convergence can be found in Appendix A. On the contrary, a join gateway requires all incoming flows to carry
enough tokens to start execution. Then, it is enough to have one Event-B event
to represent the gateway, which consumes tokens from all incoming flows.
Data. There are three kinds of data: process attributes, data stores, and activity
inputs/outputs. For each process attribute, we define a function that maps each
process instance to the runtime value of the attribute in that particular instance.
A data store is globally accessible and does not belong to a particular process.
Therefore, unlike process attributes, the data structure representing the data
store involves no process information. Finally, activities may have input and
output parameters. BPMN allows activities to have multiple sets of inputs or
outputs. However in our translation we stipulate that any flow object or subprocess has at most one input set and one output set. We also do not explicitly
represent inputs and outputs, since the runtime values of inputs/outputs are
decided by process attributes or data stores.
3.3 Events Triggers
An event either throws or catches a certain kind of triggers. The BPMN specification provides no information of trigger structures and how triggers are processed, stored, and discarded. In our understanding, each kind of triggers has its
specific processing mechanism. For instance, the trigger of a message receiving
event occurs when a desired message becomes available, and it persists until the
message is consumed. On the contrary, the trigger of a conditional event occurs
when a certain condition is fulfilled. However, if the conditional event is not
ready to be triggered, e.g, it has no incoming tokens, then the trigger immediately disappears. Based on the above discussion, we have no explicit and unified
representation for all kinds of triggers. Instead, we model the trigger behavior
implicitly in their executional contexts.
Formal Analysis of BPMN Models Using Event-B 39
3.4 Messages
Message buffers are implemented simply as sets since message order information
is absent. For each type of message, we introduce two variables to record (1)
the set of already sent messages of the type and (2) the set of messages still in
the buffer (i.e., not yet received). Note that the buffer is shared by all process
instances that may send or receive this type of messages. Sending a message
is simply to add the message into both the buffer and the set of already sent
messages, while receiving a message is to remove it from the buffer. Message
fields are defined as functions that map each message instance to the concrete
value of the corresponding field in that message.
Some message fields may contain correlation information that identifies the
intended receiver which contains matching correlation information. In the model
in Figure 1, session identifiers (sid) are used as correlation information. Each
response message contains an sid field, which can be received only by a process
instance with a matching sid as its process attribute. A request message is
used to create a new WORKER instance. Therefore, a new request message should
contain a new sid. Further detail on the translation of correlation-based message
exchanges is available in [7].
3.5 Sub-processes
In BPMN, a sub-process can be either collapsed or expanded, with the internal structure of the sub-process either hidden or revealed respectively. These
two appearances find their analogies in the refinement hierarchy of the EventB translation: The sub-process is first specified without internal detail when
the control flow of its containing process is added. The internal detail of the
sub-process is specified at later refinement levels.
For a looping sub-process, each execution creates a single outer instance, which
acts as a container for multiple inner instances. The execution of an inner
instance corresponds to a single loop iteration. Further detail on the translation
of the “collapsed view” of the loop sub-process in the FACTORY process in Figure
1 can be found in [7].
The translation of the “expanded view” is shown below. At this level we
add the outer instance attribute loop counter, and also introduce an auxiliary
variable next to control the creation of the next inner instance. In our example,
the loop condition is tested before each iteration, and therefore we initialize next
to false to enforce the checking of the loop condition before any inner instance
is created. Note that in the following code we leave out all guards and actions
inherited from abstract events.
MACHINE Level 03 Sub schedule CF
VARIABLES
......
at outer schedule loopcounter
au outer schedule next
...... EVENTS
Event act Factory schedule activate = refines act Factory schedule activate
40 J.W. Bryans and W. Wei
any
pid
child
where
... : ...... then
... : ......
act5 : au outer schedule next(child) := FALSE
act6 : at outer schedule loopcounter(child) := 0
end
Event act Factory schedule complete = refines act Factory schedule complete
any
pid
child
inners
where
... : ......
grd6 : inners = dom(outer inner schedule  {child})
grd7 : ran(inners  state inner schedule) ⊆{completed}
grd8 : at outer schedule loopcounter(child) ≥ max retry then
act1 : state outer schedule(child) := completed
act2 : tk Factory schedule gate(pid) := tk Factory schedule gate(pid)+ 1
end
Event act Factory schedule next =
any
pid
child
inners
where
... : ......
grd6 : inners = dom(outer inner schedule  {child})
grd7 : ran(inners  state inner schedule) ⊆{completed}
grd8 : at outer schedule loopcounter(child) < max retry then
act1 : au outer schedule next(child) := TRUE
end
Event evt schedule start =
any
pid
parent
outer where
... : ......
grd8 : au outer schedule next(outer)= TRUE
then
... : ......
act10 : au outer schedule next(outer) := FALSE
act11 : at outer schedule loopcounter(outer) :=at outer schedule loopcounter(outer)+1
end
END
3.6 Compensation
Compensation starts with the execution of a compensation throwing event. Each
throw event has a scope, and only activities within this scope can be compensated. An activity is within the scope of a compensation throw event if (1) the
activity is contained in the same process as the event; or (2) the event is contained in a compensation event sub-process of the process that contains the
activity.
Usually, a compensation throw event contains a reference to the activity to be
compensated. However, it is left open in the official BPMN document whether
all completed instances of the activity inside the scope will be compensated, or
only the last instance is to be compensated. In our translation, all completed
Formal Analysis of BPMN Models Using Event-B 41
instances are compensated. An activity can be compensated only after being
completed. If a compensation trigger is thrown when an activity instance is still
active, the compensation handler of the activity instance is not triggered and,
in this translation, will never be triggered unless another compensation trigger
is thrown again in the future.
The following code shows how the shipping activity in Figure 3 is compensated. au Retailer shipcomp insts records the activity instances which need
to be compensated, and au Retailer shipcomp sync is used to wait for the
completion of the involved compensations before passing tokens to outgoing
flows.
MACHINE Level 04 Retailer Data
VARIABLES
......
au Retailer shipcomp sync
au Retailer shipcomp insts
......
EVENTS
Event evt Retailer shipcomp activate = extends evt Retailer shipcomp activate
any
pid
to comp
where
grd1 : pid ∈ instances Retailer
grd2 : state Retailer(pid)= active
grd3 : tk Retailer gate shipcomp(pid) > 0
grd4 : au Retailer shipcomp sync(pid)= FALSE
grd5 : to comp ⊆ instances ship
grd6 : to comp = dom(parent ship  {pid}) ∩ dom(state ship  {completed})
then
act1 : tk Retailer gate shipcomp(pid) := tk Retailer gate shipcomp(pid) − 1
act2 : au Retailer shipcomp sync(pid) := TRUE
act3 : au Retailer shipcomp insts(pid) := to comp end
Event evt Retailer shipcomp complete = extends evt Retailer shipcomp complete
any
pid
where
grd1 : pid ∈ instances Retailer
grd2 : state Retailer(pid)= active
grd3 : au Retailer shipcomp sync(pid)= TRUE
grd4 : ran(au Retailer shipcomp insts(pid)  state ship) ⊆{compensated}
then
act1 : au Retailer shipcomp sync(pid) := FALSE
act2 : tk Retailer shipcomp chargecomp(pid):=tk Retailer shipcomp chargecomp(pid)+1
act3 : au Retailer shipcomp insts(pid) := ∅
end
Event act Retailer shipcomp = refines act Retailer shipcomp
any
pid
child
where
grd1 : pid ∈ instances Retailer
grd2 : child ∈ instances ship
grd3 : state ship(child)= completed
grd4 : parent ship(child)= pid
grd5 : child ∈ au Retailer shipcomp insts(pid)
then
act1 : state ship(child) := compensated
act2 : db order status(at Retailer order(pid)) := returned
end
END
42 J.W. Bryans and W. Wei
Fig. 3. A BPMN model for an online retailer
4 Consistency of Business Processes
We can use the Rodin toolset to examine the generated Event-B models for
properties such as deadlock and livelock. In this section we show how we may gain
further confidence in the correctness of the BPMN model by stating and proving
application-level properties as invariants within the Event-B model. We use the
online retailer model in Figure 3 as an example. The BPMN contains two extra
annotations in the top right corner. These are extra application-level consistency
conditions on the BPMN model. We anticipate these conditions to be defined
by the developer and treated by the implementor as further constraints on the
model. We show how we take account of them within the Event-B translation.
The online retailer model starts with the buyer, at which point a new instance
of the process is generated. The buyer sends a purchase order to the retailer,
which contains order and buyer information. The retailer ships the requested
item, and the buyer is then charged. If, within a specified time period, the
buyer asks to return the item, and the retailer chooses to accept the return,
then both the shipping and charging activities must be compensated – shipping
by the return of the item and charging by sending a refund to the buyer. The
consistency of the information maintained about the order status and the buyer
account must be maintained by this process.
The BPMN compensation event passes control to an associated compensating
activity (Return item and Refund in our example.) The purpose of the compensating activity is to “undo” an earlier part of the workflow. A precise specification
of the behaviour of this activity is usually left to a later stage in the development
process.
The text annotations we investigate here, such as (1) and (2) in Figure 3, give
the BPMN developer the opportunity to provide a more precise specification of
Formal Analysis of BPMN Models Using Event-B 43
required properties of this subsequent development. Text annotation (1) in states
that the order status is refunded if and only if the compensation paid is equal to
the price of the item. Translating annotation (1) extends the Event-B refinement
hierarchy with a new machine containing a new variable compensation and an
additional invariant. The variable records the compensation paid in each instance
of the retailer process. The consistency invariant introduced is formalized as
∀pid · pid ∈ instances Retailer⇒
(db order status(at Retailer order(pid)) = refunded ⇔
(compensation(pid)= price(at Retailer order(pid))))
where the order is marked as refunded only when the compensation paid is equal
to the price of the goods ordered. The event generated from the refund activity
is also extended to record the compensation paid on that order. The new event
is shown below with act4 as the additional action.
Event act Retailer chargecomp = extends act Retailer chargecomp
any
pid
child where
grd1 : pid ∈ instances Retailer
grd2 : child ∈ instances charge
grd3 : state charge(child)= completed
grd4 : parent charge(child)= pid
grd5 : child ∈ au Retailer chargecomp insts(pid) then
act1 : state charge(child) := compensated
act2 : db buyer account(at Retailer buyer(pid)) := db buyer account(at Retailer buyer(pid))+
price(at Retailer order(pid))
act3 : db order status(at Retailer order(pid)) := refunded
act4 : compensation(pid) := price(at Retailer order(pid)) end
The second annotation in Figure 3 is a property over all instances of processes.
The value in the account of any buyer should be the initial value of the account
less any purchased items. Translating annotation (2) again adds a new machine
to the model, which includes the invariant
∀b · (b ∈ BUY ERS ⇒
(db buyer account(b)= initial buyer account(b) −
Sum(ran(dom(at Buyer buyer  {b})  at Buyer order)
∩
dom(db order status  {charged, returned}))))
in which the clause ran(dom(at Buyer buyer{b})at Buyer order) identifies
all orders placed by buyer b. This is restricted to orders with status charged
or returned by the clause dom(db order status  {charged, returned}). Order
status returned identifies those orders which have been returned but not yet
refunded, and therefore still need to be included in our invariant.
Proofs. The first property results in 28 proof obligations, of which 16 are automatically discharged. The other proof obligations require expert human intervention. The second property is considerably more complex and therefore results in
44 J.W. Bryans and W. Wei
582 proof obligations, of which 300 are automatically discharged. The proving of
both properties requires the discovery and use of auxiliary invariants as lemmas.
For the second property, a total number of 88 additional invariants are added.
Currently, we need to manually discover these lemmas. However, we observe
that 30 lemmas express relations between token quantities on different sequence
flows, e.g., if the incoming flow of Charging buyer has tokens then the incoming flow of Shipping cannot have tokens. Such information can be obtained by
an automated static analysis on the control flow of the model. Therefore, it is
possible to automatically discover these 30 lemmas. In future work we will also
investigate the possibility of discovering other kinds of lemmas. Furthermore,
we observe a highly repeated pattern in the proofs that involves case splitting
to distinguish process instances. Such patterns can be implemented as proof
strategies customized for proving a certain class of invariants.
5 Enhancement of Processes Models Using Patterns
When a property is violated by a model, it is possible that the model contains
undesired behavior which can be removed by further constraining the model via
refinement steps. We may directly perform such steps on the Event-B translation
of the model in order to verify whether such refinement steps are valid, before
making changes to the original model. Moreover, refinements in Event-B can be
done automatically using patterns.
Event-B patterns [3,11,12,8] are a means of expressing reusable modeling
structures and managing effort by promoting proof re-use. In this example, we
use the type of pattern presented in [11], which provides a controlled way of
extending an Event-B development with a pre-validated refinement step. Since
the refinement step between the abstract and the concrete pattern machines has
been proved in advance, any application of the pattern results in a new, fullyproved, refinement step. The approach is automated as a plug-in for the Rodin
platform ([10]). In the example we present below, a pattern is used to correct a
previously discovered omission in a specification.
We use the shift worker scheduling process in Figure 1 as an example. The
process depicted in Figure 1 contains a timing-related fault2, which can lead to
an inconsistency in the data maintained by the business process. It is caused by
the use of the timeout at the point where worker responses are received. It arises
when a request is sent to a potential worker but no reply is received within the
allotted time. Another request is therefore sent to another available candidate.
He may accept and be assigned to the shift, after which an accept message is
received from the first worker. Now the first worker thinks he is the replacement,
but in fact the second has been chosen. We discovered this error using the Rodin
model checker ProB: we added an invariant expressing the property that at
most one worker accepts the shift at any given time. ProB found an erroneous
execution within a short time.
2 Note that the shift worker scheduling process is a simplified version of a BPMN
workflow proposal, and not a part of any real world system.
Formal Analysis of BPMN Models Using Event-B 45
Fig. 4. Structure of the timed error recovery pattern
When translated into Event-B, the flaw present in the described scenario can
be corrected using the timed error recovery pattern, shown in Figure 4 and
presented in full in [6]. It is designed to be applied to any model in which late
messages are not properly processed. When applied, a further refinement level
is added to the Event-B development. This new level contains the error recovery
behavior which ensures adequate processing of any late messages.
The concrete machine in the pattern separates normal and recovery behavior by distinguishing the receipt of messages before and after the deadline and
handling these two cases separately. Late responses are followed with a compensation event, which may be further refined depending on the way in which
recovery is implemented.
Applying the pattern requires the identification of the activities in the workflow where the timer is set and the (on time or late) replies are received. These
activities are then matched with the sending and receiving events in the pattern abstract machine (snd and rcv in the abstract machine in Figure 4). The
pattern variables must also be matched to the appropriate variables within the
development.
In the Shift Worker Scheduling model in Figure 1, the timer is set at the task
Select an available worker. The Receive response action is the point at
which messages are received. The application of the pattern introduces a new
event corresponding to rcv bad (the arrival of late replies) and given below.
Event act schedule response late =
any
pat m where
grd1 : pat m ∈ q rcv
grd2 : tt(pat m) < now then
act1 : q rcv := q rcv \{pat m}
act2 : q comp := q comp ∪{pat m}
act3 : timercvd := timercvd ∪{pat m → now} end
In this event, pat m is the message and q rcv and q comp are the messages
queued for reception and compensation respectively. The second guard requires
that the current time (now) is later that the target arrival time of the message
(tt(pat m)). On arrival, the message moves to the queue for compensation and
the time at which it is received is recorded.
The recover event refines the rcv event. As well as retaining all the functionality of rcv, it places the compensated message in the database of consistent
messages. The precise nature of the compensation activity required will vary
46 J.W. Bryans and W. Wei
according to the particular activity it is compensating, so the recover event acts
as a placeholder for a fuller description of compensation within the workflow,
which may be added (perhaps by the application of a more specific pattern) in
further refinements.
The ability to automatically add pre-validated refinement steps to generated
Event-B models can be used to support BPMN development. In our example,
the refinement step made to the Event-B translation can be re-constructed in
the original model by introducing a parallel thread to detect and react to late
messages. Such reconstruction can be achieved either through a reverse translation procedure from Event-B back to BPMN, or by building up a repository of
BPMN refinement patterns corresponding to Event-B patterns. We will explore
both possibilities in future work.
6 Related Work
Unlike our work presented in this paper, existing work in this area is largely concentrated on an examination of BPMN control flow, and does not consider data
modeling. Most of them also consider only a small fraction of the BPMN language, and put many restrictions on models that can be analyzed. [9] uses Petri
nets to formalize and analyze BPMN control flows while abstracting from data
information. The approach requires 1-safeness (i.e., having at most one token
on any sequence flow) in order to analyze exception handling for sub-processes.
On the contrary, we wish to be able to model multiple processes instances. [17]
also uses Petri nets to treat transactions and compensations, and also limits
his treatment to pure control-flow aspects of models. [20] focuses on the control
flow aspects of BPMN in a mapping to the formal workflow language YAWL.
[18] formalizes a subset of BPMN in CSP but does not consider features such
as compensations and correlation. It is also unclear how data is modeled. The
recent work to be published in [4] gives a precise and well-structured semantics
of BPMN using Abstract State Machines. This is, to our knowledge, the largest
coverage of BPMN besides ours. In [16] a subset of BPMN is translated into
the process algebra COWS in order to exploit the stochastic extensions to perform quantitative reasoning on BPMN processes. In [19] the authors present a
relative timed semantics for BPMN using CSP. This approach concentrates on
the correctness of control flow. In [14] the authors outline a framework for the
verification of business processes suggesting the use of TLA+ as well as Petri
nets.
7 Conclusions and Further Work
We have presented our work on the formal analysis of business process models through a translation into Event-B. The translation can be fully automated
and covers a large set of BPMN features. In particular, we consider the modeling
Formal Analysis of BPMN Models Using Event-B 47
of both control flow and data flow. We showed how properties can be verified
by the help of automated provers in the Rodin platform, and also showed how
Event-B patterns can be used to support the correction of design errors.
Subsequent work on this topic will be driven by our long-term goal: to allow
the BPMN developer to benefit from the improved analytic power of formal
methods (and in particular Event-B) while adding minimal extra complexity to
the design process. As an initial step, we expect to implement the presented
translation as a plug-in to the Rodin toolkit.
The two proofs of possibility presented in Sections 4 and 5 point to two different enhancements to the BPMN development method which we could aim to
support. The first, of adding annotations to BPMN, will require a definition of
the annotation language and a formalization and implementation of the rules to
translate these annotations to Event-B. The second, of using patterns to transform the generated Event-B models, would benefit from the definition of the
inverse translation from Event-B to BPMN. Note that this is not the same as a
general Event-B to BPMN translation, as we would be able to impose relatively
strong conditions on the structure (and indeed syntax) of source Event-B models
in this translation. We could also develop a library of BPMN transformations
together with their Event B patterns, and offer developers a choice from this
library in response to identified problems.
Both these approaches suffer from the high number of proof obligations which
must be manually discharged. We will therefore work on the automatic discovery
of auxiliary lemmas to assist in the proof task. Furthermore, we plan to design
and implement various proof strategies tailored for specific classes of proof obligations in order to increase the number of automatically discharged proofs.
Finally, we expect to explore the use of model-checking as a means of providing
rapid feedback to the developer on the reason for a failed proof. The challenge
here is to provide feedback in a way meaningful to the developer.
Acknowledgments
We are grateful to the anonymous reviewers for their valuable comments and to
Deploy project colleagues for their continuing fruitful collaboration. This work
has been supported by the EC FP7 Integrated Project Deploy and the EPSRC
grant TrAmS (EP/E035329/1).
References
1. Abrial, J.-R.: Modeling in Event-B: System and Software Engineering. Cambridge
University Press, Cambridge (2010)
2. Abrial, J.-R., Butler, M., Hallerstede, S., Voisin, L.: An open extensible tool environment for Event-B. In: Liu, Z., He, J. (eds.) ICFEM 2006. LNCS, vol. 4260, pp.
588–605. Springer, Heidelberg (2006)
48 J.W. Bryans and W. Wei
3. Ball, E., Butler, M.: Event-b patterns for specifying fault-tolerance in multi-agent
interaction. LNCS, vol. 5454, pp. 104–129. Springer, Heidelberg (2009)
4. B¨orger, E., S¨orensen, O.: BPMN Core Modeling Concepts. Inheritance-Based Execution Semantics. In: Handbook of database technology. Springer, Heidelberg (to
appear, 2010)
5. Bryans, J.W., Fitzgerald, J.S., Romanovsky, A., Roth, A.: Formal modelling and
analysis of business information applications with fault tolerant middleware. In:
Proc. of ICECCS 2009, pp. 68–77. IEEE Computer Society, Los Alamitos (June
2009)
6. Bryans, J.W., Fitzgerald, J.S., Romanovsky, A., Roth, A.: Patterns for modelling
time and consistency in business information systems. In: Proc. of IECCS 2010,
pp. 105–114. IEEE Computer Society, Los Alamitos (2010)
7. Bryans, J.W., Wei, W.: Formal Analysis of BPMN models using Event-B. Technical
Report CS-TR 1201, School of Computing Science, Newcastle University (May
2010)
8. Cansell, D., M´ery, D., Rehm, J.: Time constraint patterns for event B development.
In: Julliand, J., Kouchnarenko, O. (eds.) B 2007. LNCS, vol. 4355, pp. 140–154.
Springer, Heidelberg (2006)
9. Dijkman, R.M., Dumas, M., Ouyang, C.: Formal semantics and analysis of BPMN
process models using Petri nets, http://eprints.qut.edu.au/7115/01/7115.pdf
10. F¨urst, A.: Design patterns in Event-B and their tool support. Master’s thesis, ETH
Z¨urich (2009)
11. Hoang, T.S., F¨urst, A., Abrial, J.-R.: Event-B Patterns and Their Tool Support. In:
Proc. of SEFM 2009, pp. 210–219. IEEE Computer Society, Los Alamitos (2009)
12. Iliasov, A.: Design Components. PhD thesis, Newcastle University (2008)
13. K¨atker, S., Patig, S.: Model-driven development of service-oriented business application systems. In: Wirtschaftsinformatik (1), vol. 246, pp. 171–180. Osterreichische ¨
Computer Gesellschaft (2009), books@ocg.at
14. Masalagiu, C., Chin, W.-N., Andrei, S¸., Alaiba, V.: A rigorous methodology for
specification and verification of business processes. Formal Aspects of Computing 21(5), 495–510 (2009)
15. OMG. Business process model and notation (BPMN), FTF beta 1 for version 2.0,
http://www.omg.org/spec/BPMN/2.0/Beta1/PDF/
16. Prandi, D., Quaglia, P., Zannone, N.: Formal analysis of BPMN via a translation
into COWS. In: Lea, D., Zavattaro, G. (eds.) COORDINATION 2008. LNCS,
vol. 5052, pp. 249–263. Springer, Heidelberg (2008)
17. Takemura, T.: Formal semantics and verification of BPMN transaction and compensation. In: Proc. of APSCC 2008, pp. 284–290. IEEE, Los Alamitos (2008)
18. Wong, P.Y.H., Gibbons, J.: A process semantics for BPMN. In: Liu, S., Maibaum,
T., Araki, K. (eds.) ICFEM 2008. LNCS, vol. 5256, pp. 355–374. Springer, Heidelberg (2008)
19. Wong, P.Y.H., Gibbons, J.: A Relative Timed Semantics for BPMN. Electronic
Notes in Theoretical Computer Science 229(2), 59–75 (2009)
20. Ye, J.-H., Sun, S.-X., Wen, L., Song, W.: Transformation of BPMN to YAWL. In:
CSSE (2), pp. 354–359. IEEE Computer Society, Los Alamitos (2008)
Formal Analysis of BPMN Models Using Event-B 49
A Translation of Control Flows
The following shows how the end event3 in the FACTORY process in Figure 1 are
translated. This end event has two incoming flows. Every tk Factory xx xx is a
token function that maps each Factory instance to the number of tokens on the
respective flow. The end event can be executed only if the containing Factory
instance is still active. This is reflected by the guard grd2 in each Event-B event.
MACHINE Level 01 Factory CF
SEES Data Types, Processes VARIABLES
......
tk Factory schedule gate
tk Factory gate assign
tk Factory gate end
tk Factory assign end
...... INVARIANTS
... : ......
inv4 : tk Factory schedule gate ∈ instances Factory → N
... : ...... EVENTS
Event evt Factory end in1 =
any
pid where
grd1 : pid ∈ instances Factory
grd2 : state Factory(pid)= active
grd3 : tk Factory assign end(pid) > 0
then
act1 : tk Factory assign end(pid) := tk Factory assign end(pid) − 1
end
Event evt Factory end in2 =
any
pid where
grd1 : pid ∈ instances Factory
grd2 : state Factory(pid)= active
grd3 : tk Factory gate end(pid) > 0
then
act1 : tk Factory gate end(pid) := tk Factory gate end(pid) − 1
end
... ... END
3 An end event is a sink for tokens, i.e., it only consumes tokens without generating
any. Reaching an end event does not necessarily imply the completion of process
execution.